\section{Scope, Edition, and Qualifications}

\subsection{Work identity}

\begin{center}
\small
\renewcommand{\arraystretch}{1.14}
\begin{tabular}{@{}>{\raggedright\arraybackslash}p{0.24\linewidth} >{\raggedright\arraybackslash}p{0.70\linewidth}@{}}
\toprule
Rite and books & Roman Rite; \latin{Missale Romanum}, editio typica, declared typical by decree of the Sacred Congregation of Rites of 23 June 1962, operating under the \latin{Rubricae Breviarii et Missalis Romani} approved 25 July 1960 and promulgated 26 July 1960 \\
Effective law & The code binds from 1 January 1961; the Missal is the book made to fit it \\
Cutoff & 23 June 1962. No later act is stated as modifying the calendar or rubrics described here \\
Territory & The universal Roman calendar. No territorial, diocesan, religious, titular, dedication, or patronal layer is inventoried \\
Languages & English exposition; Latin from the identified witnesses; working glosses are project translations governed by the quoted Latin \\
Reader & Someone who must know what the law says and where it says it: a compiler of a calendar reference, a writer of proper guides, a student of the 1962 books \\
\bottomrule
\end{tabular}
\end{center}

\subsection{What is covered}

The complete normative temporal cycle: the six seasons and their sub-periods with their exact
boundaries; the classes and enumeration of Sundays, ferias, vigils, and octaves; the Ember and
Rogation provisions; the movable computation from Easter with both counts and their ranges; and
the resumption mechanism at \rn{18} with its Missal loci.

The complete rank and precedence machinery: the four classes and the grade conversion from the
older system; the table of precedence at \rn{91} in full; accidental occurrence and translation;
perpetual occurrence and reposition; concurrence; and commemorations with the counting rules of
\rn{111} and their Missal implementation in the orations at \rnn{433--465}.

The structure of the sanctoral: the layout of the printed calendar; the complete I class and
II class entries of the universal calendar; the rules by which particular calendars are built
(\rnn{41--58}) and by which days are assigned to new feasts (\rnn{59--62}); and the itemised
calendar changes recorded in the \latin{Variationes in Calendario}.

\subsection{What is not covered}

\begin{itemize}[leftmargin=1.3em, itemsep=0.2em]
\item An entry-by-entry inventory of the III class feasts and commemorations of the universal
calendar. Those layers are described structurally and their counts are not published.
\item Any territorial, diocesan, religious, or church calendar, and any patronal, titular, or
dedication feast. The rules governing them are given; their contents are not.
\item An Ordo for any civil year. The worked years of Section~\ref{sec:years} apply the
universal rules and are explicitly not a statement of what was or will be celebrated anywhere.
\item Mass and Office texts. Formularies are named and located; their prayers are not
reproduced, except for short rubrical sentences quoted as law.
\item The Breviary rubrics (\rnn{138--268}) and the Acta's tables of occurrence and concurrence,
which belong to the Office and are not printed in the Missal.
\item Present discipline concerning the use of the older books, which is mutable and is
deliberately not stated.
\item Any claim that the 1962 books are presently authorised in a particular place, or that any
computation here may be substituted for a competent calendar or Ordo.
\end{itemize}

\subsection{Method and source hierarchy}

The order of authority used is: the promulgating act and the promulgated code; the typical
edition of the Missal and the calendar and rubrics it prints; and, for comparison only, an
identified digitisation of a pre-1955 printing whose edition identity is unresolved. Optical
text layers were used to locate passages and never to establish a reading; every reading
reported as a divergence was read in a rendered page image. Computations were performed from the
Gregorian ecclesiastical rule and then checked against the Missal's own printed tables; where
they disagree, the disagreement is reported and not resolved by preference.

Three states are kept distinct throughout. \emph{Verified source text} is Latin read in a page
image of an identified witness. \emph{Checked paraphrase} is an English statement of a rule
checked at its exact number. \emph{Source-grounded synthesis} is an editorial conclusion drawn
from checked evidence; the Project Synthesis is labelled as such, and the assessment of the
twenty-three-Sunday readings in Section~\ref{sec:resumed} is labelled as an editorial assessment
rather than a determination.

\subsection{Ceilings on the evidence}

\begin{enumerate}[leftmargin=1.6em, itemsep=0.25em]
\item \textbf{One exemplar of each printing.} The Acta text was read in the Holy See's own
digitisation of volume 52; the Missal in the Church Music Association of America facsimile of
the typical edition. No second exemplar of either was consulted. A defect peculiar to the copy
scanned cannot be excluded from the collation findings.
\item \textbf{No official statement of the computus was located.} The paschal rule used is the
standard Gregorian ecclesiastical rule; its authority is not asserted from a Church text, and
its use here is justified only by agreement with the Missal's own printed tables over the
fifty-two dated rows checked.
\item \textbf{The pre-1955 comparison witness is a publisher's edition of unestablished printing
year.} Its title page, copyright line, and approval leaf identify it as a Benziger Brothers
\latin{editio IV iuxta typicam Vaticanam}, approved for New York on 8 December 1942 and
conformed to the Sacred Congregation's decree of 9 January 1942; the hosting item's ``1947'' is
uploader-supplied and is corroborated by nothing in the scan. It is therefore a witness to the
state of the rubrics under the 1920 Vatican typical edition, not to that typical edition's own
text, and not to the last printing before 1955. It is a Missal, so it carries neither the
Breviary rubrics nor the pre-1955 tables of occurrence and concurrence.
\item \textbf{The Missal's calendar pages were read as page images; the Acta calendar was not
collated against them.} The I class and II class lists in Section~\ref{sec:sanctoral} come from
the Missal's nine printed calendar pages read at 260 dots per inch. A machine tally of the Acta
calendar was run and was not reliable enough to publish; no count of III class entries or
commemorations is asserted.
\item \textbf{Resumed formularies partially collated.} The Third and Fourth Sundays left over
after the Epiphany were read in the facsimile and their proper Epistles and Gospels are cited;
the Fifth and Sixth were located but not separately collated, and no reading is cited for them.
\item \textbf{No Ordo was consulted.} No dated Ordo under the 1960 rubrics, for any year, was
located or read. This is why the twenty-three-Sunday case and the classification of the
Dedication of the Archbasilica remain open.
\item \textbf{Beyond the fifty-two dated rows checked, the movable dates are computed only.}
The 2026 and 2027 skeletons in Section~\ref{sec:years} were not compared with any printed
witness.
\end{enumerate}

\subsection{Rights}

The Latin quoted here is official liturgical and normative text of the Holy See; the project
claims no rights in it. Quotation is focused on the rules the reference states and does not
reproduce Mass or Office prayers. Working glosses are project-created English and are marked as
such; they are not translations of record and must not be used as though they were. The exact
digitisations relied on are recorded in the repository's source library with their retrieval
routes, hashes, and rights status: the Acta volume and the Missal facsimile are both identified
without redistribution permission established, and no payload of either is retained in the
repository. The Missal facsimile was re-downloaded during the preparation of this edition and
hashed byte-identical to the artifact record already registered.

\subsection{Review state}

Generation provenance is recorded in the terminal generation metadata. The source audit,
computation checks, and production build are recorded in the leaf's research records. No
independent liturgical, rubrical, historical, or canonical review has been performed on this
edition, and none should be inferred from the precision of its citations. This is a study aid.
It is not an official liturgical book, a critical edition, an Ordo, a canonical opinion, or a
substitute for the competent authority.
